% Chunk: Chapter 4 backmatter---Problems and References
% Source: Weinberg QFT Vol. I, physical PDF pp. 212--213 (printed pp. 189--190)
% Figures: none.
% Uncertainties: none.

\chapterexercisehook{04}

\chapterbackmatter{References}

\begin{enumerate}
\item \label{ref:4.1}The cluster decomposition principle seems to have been first stated
explicitly in quantum field theory by E. H. Wichmann and J. H. Crichton,
\emph{Phys. Rev.} \textbf{132}, 2788 (1963).

\item \label{ref:4.2}See, e.g., B. Bakamjian and L. H. Thomas, \emph{Phys. Rev.}
\textbf{92}, 1300 (1953).

\item \label{ref:4.3}For references, see T. D. Lee and C. N. Yang, \emph{Phys. Rev.}
\textbf{113}, 1165 (1959).

\item \label{ref:4.4}J. Goldstone, \emph{Proc. Roy. Soc. London} A\textbf{239}, 267 (1957).

\item \label{ref:4.5}N. M. Hugenholtz, \emph{Physica} \textbf{23}, 481 (1957).

\item \label{ref:4.6}See, e.g., R. Kubo, \emph{J. Math. Phys.} \textbf{4}, 174 (1963).

\item \label{ref:4.7}E. C. Titchmarsh, \emph{Introduction to the Theory of Fourier
Integrals} (Oxford University Press, Oxford, 1937): Section 1.8.

\item \label{ref:4.8}L. D. Faddeev, \emph{Zh. Eksper. i Teor. Fiz.} \textbf{39}, 1459
(1961) (translation: \emph{Soviet Phys.---JETP} \textbf{12}, 1014 (1961));
\emph{Dokl. Akad. Nauk. SSSR} \textbf{138}, 565 (1961) and \textbf{145}, 30
(1962) (translations \emph{Soviet Physics---Doklady} \textbf{6}, 384 (1961)
and \textbf{7}, 600 (1963)).

\item \label{ref:4.9}S. Weinberg, \emph{Phys. Rev.} \textbf{133}, B232 (1964).
\end{enumerate}
